\scanpage{11}
\mspage{9}

cette opération, via l'isomorphisme (2.4.2), \deleted{n'est} n'est autre que celle qui provient de \deleted{l'isomorphisme de} \inserted{l'opération naturelle} de $\pi_0$ sur les facteurs $\mathbb Z_\ell(1)$, via l'expression (2.3.3) de $\mathbb Z_\ell(1)$.

Bien entendu, pour tout $\ell$ \unsure{\deleted{\unsure{on a}}\inserted{il existe}|\uline{il existe}} \unsure{des} isomorphismes
\[
\mathbb Z_\ell(1)\simeq\mathbb Z_\ell,
\]
dont le choix revient au choix d'un syst. cohérent de racines \inserted{primitives} $\ell^\nu$-ièmes de 1 dans $\bar K$; Mais on fera attention qu'il n'y a pas de façon naturelle de faire ce choix, et que c'est l'expression intrinsèque (2.4.2) qui seule \deleted{\unsure{peut}} tient compte des opérations naturelles de $\pi_0$ sur le premier membre.
\margin{10}
